\documentclass{article}
\usepackage[utf8]{inputenc}
\usepackage[T1]{fontenc}
\usepackage[margin=2cm]{geometry}
\usepackage{hyperref}
\sloppy
\begin{document}
\section*{Historical source}
Encyclopaedia Britannica scan, source row 258. Parse 3.5 extracted the text. Super 3.5 VL corrected it from an enlarged crop of each block. It can still contain recognition errors.
Page headers, footers, and page numbers are removed.

\subsection*{Untitled continuation (E0012)}
\textit{Continues an entry from an earlier page.}

lordship, and some time after the great feel. Edward being returned to England, after having carried on a very successful war in France, found his exchequer in great disorder. The duke of Lancaster, one of his sons, at the head of several lords, having brought complaints against the clergy, who then enjoyed most posts in the kingdom, the king removed them from their employments. But the laymen, who were raised to them, behaved so ill, that the king was forced to restore the ecclesiastics. The duke of Lancaster showed strong animosity to the clergy, and set every engine at work to ruin Wykeham. He impeached him of extortion, and of disguising things, and obliged him to appear at the King's-bench. He got such judges appointed as condemned him; and not satisfied with depriving him of all the temporalities of his bishopric, he advised Edward to banish him: but this prince rejected the proposal, and afterward restored to Wykeham all that he had been divested of. Richard II. was but eleven years old when Edward died; whereby the duke of Lancaster had an easy opportunity of reviving the accusations against the bishop of Winchester: nevertheless Wykeham cleared himself. Then he founded two noble colleges, the one in Oxford, the other in Winchester. Whilst he was exerting his utmost endeavours to improve these two fine foundations, he was recalled to court, and in a manner forced to accept of the office of lord high-chancellor in 1389.—Having excellently discharged the duties of that employment for three years, he obtained leave to resign it, foreseeing the disturbances that were going to break out. Being returned to his church, he finished his college, and built there so magnificent a cathedral, that it almost equals that of St. Paul's in London. He laid out several funds in things advantageous to the public and to the poor; notwithstanding which, in 1397 he was in great danger; for he and some others were impeached of high-treason in open parliament: however, he was again fully cleared. From that time till his death he kept quiet in his diocese, and there employed himself in all the duties of a good prelate. He died in 1404, in the 81st year of his age.

\subsection*{WILLIAM (E0013)}
WILLIAM, the name of several kings of England. See ENGLAND, n 87-92, and BRITAIN, n 302.

Fort-William, a fortress in the Highlands of Scotland, erected in king William's reign, as was also a small town adjoining, called Maryburgh, in honour of his queen. It is situated in Invernesshire, on a narrow arm of the sea called Loch Eil, which might easily, by a very short canal, be united to the Western sea. Fort-William is of a triangular form, having two bastions, and is capable of admitting a garrison of 800 men; but could not be defended against an attack, as it is commanded by several hills in the neighbourhood.

William's Fort, is a factory of Asia belonging to the East-India company, seated on one of the branches of the river Ganges, in the kingdom of Bengal. The fort was first built in the shape of an irregular tetragon of brick and mortar; and the town has nothing regular in it, because every one built a house as he liked best, and for his own convenience. The governor's house is within the fort, and is the best piece of architecture in those parts. Here there are also convenient lodgings for the factors and writers, with store-houses for the company's goods, and magazines for ammunition. About 50 yards from the fort is the church, built by the charity of merchants residing here. The town is called Calcutta, and has a pretty good hospital for the sick, though few come out of it alive. It is governed by a mayor and aldermen, as most of the company's factories in the East Indies now are. In 1757 it was fur-

prised by the nabob of Bengal, who took it, and put most of those that had made resistance into a place called the Black Hole, where most of them were smothered. This nabob was afterwards killed, and another set up in his room, more friendly to the English; and the factory was re-established. K. Long. 86. o. N. Lat. 22. 27.

\subsection*{WILLIAMSBURG (E0014)}
WILLIAMSBURG, a town of North America, in Virginia, and formerly capital of that state. It is situated between two creeks; one falling into James, and the other into York River. The distance of each landing place is about a mile from the town, which, with the disadvantage of not being able to bring up large vessels, and the want of enterprise in the inhabitants, has occasioned its decay. Here is a college, designed for the education of the Indians, but which, on account of their aversion to learning, never answered the purpose. It is 60 miles east of Richmond. W. Long. 76. 30. N. Lat. 37. 10.

\subsection*{WILLIAMSTADT (E0015)}
WILLIAMSTADT, a sea-port town of Holland. It is a handsome strong place, and the harbour is well frequented. It was built by William prince of Orange in 1585, and in 1732 belonged to the stadtholder of Friesland. The river near which it is built is called Butterdiep or Holland Diep, and is one of the bulwarks of the Dutch on the side of Brabant, where they always keep a garrison. This place made a gallant defence in 1793 against the French, who were obliged to raise the siege. It is 15 miles north-east of Bergen-op-Zoom, and 12 south-west of Dort. E. Long. 4. 30. N. Lat. 51. 39.

WILLIS (Dr Thomas), a celebrated English physician, was born at Great Bodwin, in Wiltshire, in 1621, and studied at Christ-church college, Oxford. When that city was garrisoned for the king, he, among other scholars, bore arms for his Majesty, and devoted his leisure hours to the study of physic. The garrison of Oxford at length surrendering to the parliament, he applied himself to the practice of his profession; and soon rendered himself famous by his care and skill. He appropriated a room as an oratory for divine service according to the church of England, whither most of the loyalists in Oxford daily resorted. In 1660, he became Sedleian professor of natural philosophy, and the same year took the degree of doctor of physic. In 1664, he discovered the famous medicinal spring at Althorp, near Brackley. He was one of the first members of the Royal Society, and soon made his name illustrious by his excellent writings. In 1666, after the fire of London, he removed to Westminster; and his practice became greater than that of any of the physicians his contemporaries. Soon after his settlement in London, his only son Thomas falling into a consumption, he sent him to Montpellier in France for the recovery of his health; and it proved successful. His wife also labouring under the same disorder, he offered to leave the town; but she, not suffering him to neglect the means of providing for his family, died in 1670. He died at his house in St Martin's in 1675, and was buried near her in Westminster-abbey. Dr Willis was extremely modest and unambitious, and refused the honour of knighthood. He was remarkably pious: As he rose early in the morning, that he might be present at divine service, which he constantly frequented before he visited his patients, he procured prayers to be read beyond the accustomed times while he lived; and at his death settled a stipend of 20l. per annum to continue them. He was a liberal benefactor to the poor wherever he came, having from his early practice allotted part of his profits to charitable uses. He was exact and regular in all his hours; and though his table was the resort of most of the great men of London, yet he was remark-


\section*{Source record}
\begin{description}
\item[Scan ID] \texttt{britannica:309f765b1ecc83e0a7f7762e44b4a927a7555443:train:258}
\item[Image file] \texttt{page-0258-631a7ae7df34.jpg}
\item[Image SHA-256] \texttt{631a7ae7df340a64a33318be0161e911b328ec2b7354754da2dc37a1606ece4c}
\item[Dataset] NationalLibraryOfScotland/encyclopaedia\_britannica\_illustrated at revision \texttt{309f765b1ecc83e0a7f7762e44b4a927a7555443} (CC0-1.0)
\item[Extraction] model \texttt{nvidia/nemotron-3.5-parse}, raw reply \texttt{extract/raw/page-0258-631a7ae7df34.json}
\item[Raw reply SHA-256] \texttt{d18230b8028dd3fcc1ada71c66a6a2bf82f9d5804ee6475bc4ef9975e75e7c4f}
\item[Correction] model \texttt{nvidia/Nemotron-3.5-Super-VL}, block crops, raw replies \texttt{correct/raw/crop\_nothink/page-0258-631a7ae7df34-b*.json}; super\_vl: 8
\item[Running header] none
\end{description}
\end{document}
